% STATUS: written 2026-11-06 (English gate clean; sources read in full as recorded below).
% Sources read in full: Raudenbush 1934 (Trans. AMS 36, 361--368).
% Not yet read: Ritt 1932 (AMS Colloquium Publ. 14), Baer 1927; described
% here only as Raudenbush reports them. Rewrite from Ritt directly once read.
\section[Differential rings and fields: Ritt and the New York school][Differential rings and fields]{Differential rings and fields: Ritt and the New York school}\label{sec:origins:differential-algebra}

Differential algebra began as an algebraic theory of systems of algebraic
differential equations. It started with the Colloquium volume of J.~F.~Ritt of
1932, \emph{Differential equations from the algebraic standpoint}. Its first
abstract definitions appear in a paper of Raudenbush of
1934~\cite{Raudenbush1934}. Raudenbush was a member of the school of Ritt.

Raudenbush begins with the problem. He reports two results of Ritt. First, Ritt had introduced the idea
of an irreducible system of algebraic differential equations. Second, he had
shown that every system is equivalent to a finite set of irreducible systems.
Raudenbush wants an abstract ideal theory. The theorem of Ritt should then
follow from it~\cite[p.~361]{Raudenbush1934}.

The objects of the theory are polynomials in unknowns $y_1, \ldots, y_n$ and their
derivatives. Following Ritt, Raudenbush calls them ``forms''. Their
coefficients lie in a differential field.

\begin{definition}[{\cite[p.~361]{Raudenbush1934}}]\label{def:origins:differential-field}
A \emph{differential field} is a field $K$ with a map $\partial\colon K \to K$
such that
\[
\partial(a+b) = \partial a + \partial b, \qquad \partial(ab) = \partial(a)\, b + a\, \partial(b)
\qquad (a, b \in K).
\]
The polynomials in unknowns $y_1, \ldots, y_n$ and all their derivatives, with
coefficients in $K$, form a \emph{differential ring}, and a
\emph{differential ideal} is an ideal $I$ of it with $\partial(I) \subseteq I$.
\end{definition}

Raudenbush writes $a_1$ for the derivative of $a$. We write $\partial a$.

For abstract differential fields, Raudenbush refers to a work of R.~Baer. Its
title is \emph{Algebraische Theorie der differentiierbaren
Funktionenk\"orper}. It appeared in the proceedings of the Heidelberg
Academy~\cite[p.~361, footnote]{Raudenbush1934}.

He also records an example of Ritt. In this example a differential ideal has
no finite basis in the ordinary sense. For this reason the paper works with a
special class of ideals~\cite[p.~361]{Raudenbush1934}.

The theory that grew from here is the subject of
Chapter~\ref{ch:differential-algebra}. It includes the later work of Kolchin.

With Ritt, the first part of our story ends. All the rings in it are
commutative. But in 1925 physicists had already met the product rule on
matrices, where the order of the factors matters.
